% !TeX root = main.tex
\chapter*{ABSTRACT}
\phantomsection
\addcontentsline{toc}{chapter}{ABSTRACT}

Denysiuk D.~O. Information technology for detecting signs of anomalous content in multimedia data within a web environment. Qualifying scientific work presented as a manuscript.

Dissertation for the degree of Doctor of Philosophy in specialty 126 ``Information Systems and Technologies'', field of knowledge 12 ``Information Technologies''. Khmelnytskyi National University, Khmelnytskyi, 2026.

The dissertation addresses the scientific and applied task of forming a justified assessment of signs of anomalous content in multimedia data from incomplete and heterogeneous observations obtained during processing in a web system. Formal compliance with a declared format and visual plausibility do not exclude structural, spectral and visual, or event and behavioural deviations. Analysing a single representation in isolation does not establish a traceable relationship between carrier properties, a processing operation, and a possible consequence. In this dissertation, anomalous content denotes data or structural values that deviate from established format constraints, statistical descriptions, or a controlled processing route. An observed deviation is not equated with a proven threat, executable code, or malicious code.

The research aims to develop an information technology for detecting specified structural, spectral, visual, event, and behavioural signs of anomalous content in multimedia data within a web environment, based on a multimodal hidden-threat representation model, a method for analysing spectral, visual, and structural features, and a method for behavioural alignment of features of different origin. The research object is the process of forming a decision on signs of anomalous content while multimedia data are received and processed by web-system components. The research subject comprises the multimodal hidden-threat representation model, the method for analysing spectral, visual, and structural features, the method for behavioural alignment of features of different origin, and the information technology for detecting signs of anomalous content in multimedia data within a web environment.

The study employs systems analysis, mathematical modelling, digital image and video processing, structural analysis of multimedia containers, machine-learning methods, descriptive statistics, grouped resampling, and paired comparison of configurations. Systems analysis and mathematical modelling were used to formalise the object route, multimodal representation, and decision rule, while digital processing and structural analysis formed the static features. Statistical methods supported the assessment of confusion matrices, indicator uncertainty, and differences between configurations.

The first contribution is a multimodal hidden-threat representation model in which a multimedia object is treated as a multilevel information container. The immutable byte representation, controlled decoded content, structural map, technical context, and causally linked event log remain separate entities connected by a common identifier and provenance data. Structural features, spectral and visual features, and features of events and behaviour are formally linked in representation $Z$ while preserving partial scores, localisations, and an availability mask. This structure separates risk assessment, confidence, analytical state, and technological action and prevents an unavailable feature group from being interpreted as evidence of normality. The complete representation $Z$ with causally linked events was not formed experimentally.

The method for analysing spectral, visual, and structural features was improved. Its spectral and visual branch forms high-frequency residuals, discrete cosine and discrete wavelet transform characteristics, entropy features, and temporal frame aggregates. The structural branch inspects headers, container-element boundaries, service fields, termination markers, and additional byte regions. Separate scores and localisations from the two branches are retained until their integration. Reduced numerical projections were initially evaluated for PNG and MP4; transfer of fixed parameters to JPEG, WebP, and WebM was assessed without retraining or threshold adjustment. The complete static description and its localisations were not evaluated experimentally.

A method for behavioural alignment of features of different origin was developed. It matches a static description with processing events of the same object according to identifier, temporal boundaries, phase, and causal relationship. The method defines input and output contracts, causal-phase admission, event encoding, controlled integration, boundary states, and verification conditions. A reduced format-aware model for aligning static scores was implemented for experimental evaluation. It processes scores already formed by the first method and is not an implementation of the complete encoder of events and behavioural features or the cross-alignment component.

The information technology was further developed by organising the model and the two methods into five stages: preprocessing, feature formation, feature integration, decision-making, and result registration. The original object remains immutable, a reconstructed object enters a separate processing cycle, and the typed output contains component versions, localisations, an availability mask, and a log reference. The analytical result is separated from the policy-defined action of allowing, expert review, reconstruction, or blocking. State transitions $S_1$--$S_5$, atomic registration, repeated execution, recovery, and termination without an external action after two classes of technical failure were verified locally. The practical value of the results lies in specifying a traceable technological process between the multimedia input channel and the application logic of a web system, together with a machine-readable record for repeated verification.

The experimental study comprised six series: one retrospective, one internal group-separated, and four external series, supplemented by a separate resource evaluation and a local technology-state audit. According to the preserved final-run summary, the internal series contained $200$ primary PNG and MP4 groups and $400$ derived objects split in a $60:20:20$ ratio. On the held-out subset of $80$ objects, the fused numerical projection $Z_c$ correctly identified $40$ control objects and $38$ modified objects, with no false positives and two false negatives. The confusion matrix independently yields an accuracy of $0.9750$, a recall of $0.9500$, and $F_1=0.9744$; the reported area under the receiver operating characteristic curve of $0.9869$ cannot be independently recomputed because continuous object-level scores are unavailable.

The four external datasets contained $38$ natural video recordings from $12$ devices and $24$ provenance groups. The first external series of $56$ derived MP4 objects did not confirm an advantage of the fused numerical projection over the structural configuration: their balanced accuracies were $0.7857$ and $0.8333$, respectively. In the multiformat series, $500$ base clips produced $10\,000$ derived objects in PNG, JPEG, WebP, MP4, and WebM. Balanced accuracy was $0.7636$ for the structural, $0.5001$ for the spectral and visual, and $0.7916$ for the fused configuration. The number of derived objects was not treated as the number of independent sources.

The reduced static-score alignment model was evaluated separately on $2\,000$ objects formed from $100$ base clips of eight new recordings captured by four devices. After fixing the model and threshold, balanced accuracy changed from $0.7827$ for the fused configuration to $0.8307$ for the alignment model. The model removed $76$ false-positive decisions but introduced $84$ additional false negatives and did not detect $500$ objects containing only the spectral modification. The exact two-sided McNemar test for ordinary accuracy yielded $p=0.5801$. The result therefore indicates a change in the balance between correct identification of control objects and recall rather than unconditional superiority of the alignment model.

The retrospective assessment of $110$ objects confirmed the arithmetic consistency of the aggregate confusion matrix and its derived indicators; the area under the receiver operating characteristic curve was not determined because continuous object-level scores were unavailable. The four recorded operations averaged $418.09$~ms, but this value does not represent the complete five-stage process. Across five resource runs, the 95th-percentile times for the structural, fused, and alignment configurations were $2.136$, $853.971$, and $1078.904$~ms. In the local state audit, $35$ cycles produced complete $S_5$ records, whereas ten expected failures terminated without $S_5$ and without an external action; one four-process run achieved $22.827$ objects per second. The integral rating $K_I$ was used only as a secondary means of comparing detection validity, processing time, memory use, and stability across formats.

The conclusions are limited by the dependence of derived objects on common base clips and by the number of independent devices. Non-destructive feature re-extraction from the same $56$ and $2\,000$ MP4 files preserved the structural branch's balanced accuracy of $0.8333$ but changed the fused projection from $0.7857$ to $0.6905$ and from $0.7987$ to $0.6520$; the current reference spectral extractor is therefore not treated as numerically identical to the missing historical implementation. Among $1\,500$ benign structurally varied controls, the structural false-positive rate for private extensions or padding was $1.0000$ in every format; the structural output is therefore not presented as self-sufficient evidence of concealed influence. In an exploratory paired analysis, each of the $2\,500$ format representations containing an introduced spectral modification retained a non-zero difference after decoding; however, the original spatiotemporal localisation could not be assessed because the masks formed before encoding were unavailable. Causally linked event logs and the complete representation of events and behaviour were not collected, and the available technical feature did not trigger in $4\,832$ physically separate derived video objects, which were not treated as independent sources. The complete behavioural alignment method, detection of executable malicious code, reconstruction, execution of external policy actions, and distributed or sustained industrial load were not evaluated experimentally. The results therefore support the static basis of the first method and the local prototype implementation of the state transitions only within the evaluated data and do not justify unconditional automated blocking.

Nine scientific works have been published on the dissertation topic. Two articles in Ukrainian professional scientific journals report the main results, and seven publications document their approbation.

\textbf{Keywords:} anomalous content, multimedia data, structural analysis, analysis of spectral and visual features, multimodal representation, information technology, web environment.

\clearpage
